\documentclass[10pt,a4paper]{amsart}
\usepackage[margin=19mm]{geometry}
\usepackage{amsmath,amssymb,mathtools}
\usepackage[scr=rsfs,cal=euler]{mathalfa}
\usepackage{xfrac}
\usepackage[unicode,hidelinks,bookmarks=false]{hyperref}
\newcommand{\ie}{i.e.\ }
\newcommand{\cf}{cf.\ }
\newcommand{\resp}{respectively\ }
\newcommand{\Subsection}{\subsection}
\providecommand{\nobreakdash}{\nobreak\hbox{-}}
\setlength{\emergencystretch}{2em}
\multlinegap0pt
\usepackage{amssymb}
\usepackage{xcolor}
\newtheorem{remark}{Remark}
\newtheorem{lemma}{Lemma}
\def\N{{\mathbb N}}
\def\Q{{\mathbb Q}}
\def\Z{{\mathbb Z}}
\def\ln{\operatorname{ln}}
\def\ord{\operatorname{ord}}
\title{The Reidemeister and the Nielsen numbers: growth rate, asymptotic behavior, dynamical zeta functions and the Gauss congruences}
\author{Alexander Fel'shtyn, Mateusz Slomiany}
\date{Source: arXiv:2603.08527v1 (CC BY 4.0)}
\begin{document}
\maketitle

\noindent\emph{Source and licence.} The Reidemeister and the Nielsen numbers: growth rate, asymptotic behavior, dynamical zeta functions and the Gauss congruences. Version arXiv:2603.08527v1 (CC BY 4.0), distributed by arXiv under the Creative Commons Attribution 4.0 International licence, http://creativecommons.org/licenses/by/4.0/, as recorded in the arXiv licence metadata for this exact version and shown on the abstract page https://arxiv.org/abs/2603.08527v1. Full licence notice: \texttt{licenses/arxiv-2603.08527-CC-BY-4.0.txt}.\par\smallskip

\noindent\textit{Editorial scope.} Counted unit: the article's lemma thm:padic-log-geq-1/n (source0.tex lines 507--514). The article's complete original proof of it is delivered with it (source0.tex lines 515--576, one \texttt{proof} environment of 62 lines). No mathematics is written, repaired or re-derived: the statement and the proof are the article's own lines, in the article's own notation, and no retained line is substituted or re-flowed.  Retained uncounted as the source-local context the proof needs: lines 466--468; lines 481--484.  Nothing was omitted from any retained span, so the package carries no omission record.  No retained line contains a citation, so the wrapper carries no bibliography and no bibliography entry.  No retained line contains a figure environment, an image-inclusion command or drawing code, so no image is delivered and the package declares no asset dependency.  Editorial formatting only, in addition: an AMS \texttt{amsart} preamble that restates the article's own declarations and macros used by the retained lines, namely the environments \texttt{remark}, \texttt{lemma} on separate counters, together with the standard packages those lines need.  The retained environments print in this document as Remark 1, Lemma 1 in order of appearance, whereas the article prints the same items with the numbering of its own full document.  The complete native source of the article as distributed in the arXiv e-print for this exact version is bundled unchanged as \texttt{sources/196011/source0.tex}.
\begin{equation}\label{eq:padic-triangle}
  \big( |x|_p < |y|_p \big) \quad \Longrightarrow \quad |x-y|_p = |y|_p.
\end{equation}

\begin{remark}\label{thm:possible-values-of-padic-norm}
  Possible values for $|\cdot|_p$ on $\Q_p$ are in the set $\{p^n\}_{n \in \Z} \cup \{0\}$.
  In particular the greatest possible value $<1$ is $1/p$.
\end{remark}

\begin{lemma}\label{thm:padic-log-geq-1/n}
  Let $\xi$ be an algebraic number over $\Q_p$ such that $|\xi|_p = 1$ but $\xi$ is not a root of unity.
  Then there is a constant $C$ such that for every $n \in \N$ the inequality
  \begin{equation}\label{eq:1/n-leq-x-1}
    \frac{1}{n} \leq C \cdot \left| \xi^n - 1 \right|_p
  \end{equation}
  holds. The constant does not depend on $n$.
\end{lemma}

\begin{proof}
  For $|1|_p = 1 = |\xi|_p$, from (\ref{eq:padic-triangle}) we have $|\xi^n-1|_p \leq 1$ for all $n \in \N$.
  If for some $n$ there is equality $|\xi^n-1|_p = 1$ then (\ref{eq:1/n-leq-x-1}) holds for $C=1$

  Now consider $n$ for which $|\xi^n-1|_p < 1$.
  Denote the degree $\left[ \Q_p(\xi) : \Q_p \right] =: k$.
  From remark \ref{thm:possible-values-of-padic-norm} and
  the~definition of the $p$-adic norm for algebraic numbers we have
  \begin{equation}\label{eq:bound-for-padic-norm-of-x-1}
      |\xi^n-1|_p \leq \frac{1}{\sqrt[k]{p}}.
  \end{equation}
  Therefore we can compute the $p$-adic logarithm and
  \begin{align*}
    |n|_p \cdot |\log_p (\xi)|_p &= |\log_p (\xi^n)|_p \\
    &= \left|(\xi^n - 1) - \frac{(\xi^n - 1)^2}{2} + \frac{(\xi^n - 1)^3}{3} - \dots \right|_p \\
    &\leq \max_m \left( \left|\frac{(\xi^n - 1)^m}{m} \right|_p \right) \\
    % &= |\xi^n - 1|_p.
  \end{align*}
  Note that using (\ref{eq:bound-for-padic-norm-of-x-1}) we get
  \begin{equation}\label{eq:padic-maximum-term}
    0 \leq \left| \frac{(\xi^n - 1)^m}{m} \right|_p
    = p^{\ord_p m} \cdot \left| \xi^n - 1 \right|_p^m
    \leq m \cdot \left| \xi^n - 1 \right|_p^m
    \leq m \cdot \left( \frac{1}{\sqrt[k]{p}} \right)^m.
  \end{equation}
  For a constant $q \in (0, 1/\sqrt[k]{p}]$ the derivative
  $\frac{\partial}{\partial x} xq^x = q^x(1+x\ln q)$ vanishes only for
  $x = -1/\ln q$, is positive for smaller $x$ and negative for greater ones.
  Note that $-1/\ln q$ is an increasing function therefore in our situation $-1/\ln q \ \leq \ k/\ln p$.
  It is possible then that for a finite number of initial $m$s the terms over which we take the maximum
  actually grow, but eventually there is a maximum $M = M(n)$ and the sequence goes to $0$ as $m \to \infty$.

  In other words there exists $m_0 = m_0(n)$ satisfying $1 \leq m_0 \leq k/\ln p$, for which
  \begin{equation*} %\label{eq:max-as-|xn-1|-times-sth}
    \max_m \left( \left|\frac{(\xi^n - 1)^m}{m} \right|_p \right)
    = |\xi^n-1|_p \cdot \frac{\left| \xi^n-1 \right|_p^{m_0-1}}{|m_0|_p}
  \end{equation*}
  We can bound from above the right side similarly as above:
  \begin{align*}
    |\xi^n-1|_p \cdot \frac{\left| \xi^n-1 \right|_p^{m_0-1}}{|m_0|_p}
    &=  |\xi^n-1|_p \cdot \Big( p^{\ord_p m_0} \cdot \left| \xi^n-1 \right|_p^{m_0-1} \Big) \\
    &\leq |\xi^n-1|_p \cdot \Big( m_0 \cdot \left| \xi^n-1 \right|_p^{m_0-1} \Big)
  \end{align*}
  For $x \in [1, k/\ln p]$ the derivative $\frac{\partial}{\partial q} xq^{x-1} = x(x-1)q^{x-2} \geq 0$ so
  \begin{align*}
    |\xi^n-1|_p \cdot \Big( m_0 \cdot \left| \xi^n-1 \right|_p^{m_0-1} \Big)
    \ \leq \  |\xi^n-1|_p \cdot \left( m_0 \cdot \frac{1}{(\sqrt[k]{p})^{m_0-1}} \right)
  \end{align*}
  Note that
  \[
    M^* := \max_{1 \ \leq \ m \ \leq \ k/\ln p} \left\{ m \cdot p^{-(m-1)/k} \right\}
    \ \geq \ m_0 \cdot p^{-(m_0-1)/k}
  \]
  for all possible $m_0$, therefore $M^*$ does not depend on $n$. We got
  \[
    \frac{1}{n} \ \leq \ |n|_p \ \leq \ |\xi^n-1|_p \cdot \frac{M^*}{|\log_p(\xi)|_p}
  \]
  and (\ref{eq:1/n-leq-x-1}) holds for $C = M^* / |\log_p(\xi)|_p$.

  In general (\ref{eq:1/n-leq-x-1}) holds for all $n \in \N$ regardless of actual values of $|\xi^n-1|_p$
  if we put $C := \max \left\{1, M^* / |\log_p(\xi)|_p \right\}$.
\end{proof}

\end{document}
